\section*{Hoofdstuk \ref{ch:g12:selection-drift-speciation} --- Selectie, drift en soortvorming}

\begin{solution}{exo:g12:selection-drift-speciation:1}
Het deel van alle kopieën van een \omterm{def:g10:universal-dna:gene}{gen} in een \omterm{def:g12:selection-drift-speciation:population}{populatie} dat een bepaald
\omterm{def:g10:universal-dna:gene}{allel} is. Met $p = 0.8$, $q = 0.2$: $AA$ 64\%, $Aa$ 32\%, $aa$ 4\%.
\end{solution}

\begin{solution}{exo:g12:selection-drift-speciation:2}
Selectie: een verschil in voortplanting tussen dragers van verschillende
\omterm{def:g10:universal-dna:gene}{allelen}, door de omgeving bepaald, dat de \omterm{def:g12:selection-drift-speciation:population}{frequenties} in een richting
verandert. Drift: een willekeurige schommeling van de \omterm{def:g12:selection-drift-speciation:population}{frequenties} door de
greep uit een eindige \omterm{def:g12:selection-drift-speciation:population}{populatie}, zonder richting. De selectie hangt af van
wat het \omterm{def:g10:universal-dna:gene}{allel} doet; de drift niet.
\end{solution}

\begin{solution}{exo:g12:selection-drift-speciation:3}
Ongeveer 75\% in 1880, toen roet de stammen zwart had gemaakt; ongeveer
30\% in 1980, nadat wetten op schone lucht de stammen weer bleek hadden
laten worden.
\end{solution}

\begin{solution}{exo:g12:selection-drift-speciation:4}
De \omterm{def:g10:universal-dna:gene}{allelen} van de volgende generatie zijn een greep uit die van de vorige;
een kleine greep wijkt sterker van de bron af dan een grote, dus
schommelen de \omterm{def:g12:selection-drift-speciation:population}{frequenties} meer en kan een \omterm{def:g10:universal-dna:gene}{allel} door toeval verdwijnen.
\end{solution}

\begin{solution}{exo:g12:selection-drift-speciation:5}
Een onderbreking van de genenuitwisseling tussen twee \omterm{def:g12:selection-drift-speciation:population}{populaties}, lang
genoeg om hen te laten uiteengroeien tot zij zich niet meer onderling
kunnen voortplanten.
\end{solution}

\begin{solution}{exo:g12:selection-drift-speciation:6}
Kopieën van $a$: $160 + 40 = 200$ op 1000, dus $q = 0.2$, $p = 0.8$.
Verwacht: $0.64 \times 500 = 320$, $0.32 \times 500 = 160$, $0.04
\times 500 = 20$ --- precies de telling: de \omterm{def:g12:selection-drift-speciation:population}{populatie} verkeert in de
referentietoestand.
\end{solution}

\begin{solution}{exo:g12:selection-drift-speciation:7}
$q^2 = 0.0001$: één op tienduizend vertoont het kenmerk; $2pq \approx
0.02$: één op vijftig draagt het. Negenennegentig procent van de kopieën
zit in heterozygoten, op wie selectie tegen het kenmerk geen vat heeft.
\end{solution}

\begin{solution}{exo:g12:selection-drift-speciation:8}
Beide bereikten in generatie 12 een uiteinde, de een vastgelegd op 100\%,
de ander verloren. De derde dwaalde in 20 generaties zonder een van beide
grenzen te raken --- opnieuw het toeval; met meer tijd zou zij dat wel
doen.
\end{solution}

\begin{solution}{exo:g12:selection-drift-speciation:9}
Bleke vlinders bleven zich voortplanten op het onvervuilde platteland,
waar zij in het voordeel waren, en vlinders vliegen: elke generatie
brachten trekkers het bleke \omterm{def:g10:universal-dna:gene}{allel} weer de beroete streek in, zodat het
nooit verdween.
\end{solution}

\begin{solution}{exo:g12:selection-drift-speciation:10}
Twintig dieren dragen hoogstens veertig kopieën van elk \omterm{def:g10:universal-dna:gene}{gen}, een
piepkleine greep uit de oorspronkelijke \omterm{def:g10:universal-dna:gene}{allelen}; de meeste gingen in de
flessenhals verloren, en de drift in de kleine \omterm{def:g12:selection-drift-speciation:population}{populatie} die volgde nam er
nog meer weg. De aantallen kwamen terug; de \omterm{def:g10:universal-dna:gene}{allelen} kunnen dat niet,
behalve door nieuwe \omterm{def:g11:mutations:mutation}{mutatie}.
\end{solution}

\begin{solution}{exo:g12:selection-drift-speciation:11}
Tien \omterm{def:g12:selection-drift-speciation:population}{populaties} die dezelfde kant op bewegen, na dezelfde verandering in
de omgeving, met een aannemelijk mechanisme: selectie. Eén piepkleine
\omterm{def:g12:selection-drift-speciation:population}{populatie} van dertig, zonder gemelde oorzaak: drift is de
standaardverklaring tot er een mechanisme is aangetoond.
\end{solution}

\begin{solution}{exo:g12:selection-drift-speciation:12}
$AA$ $0.85^2 \approx 72\%$, $Aa$ $2 \times 0.85 \times 0.15 \approx
26\%$, $aa$ $0.15^2 \approx 2\%$. De $aa$ sterven en nemen kopieën van
$a$ mee; maar de $Aa$, een kwart van de \omterm{def:g12:selection-drift-speciation:population}{populatie}, overleven malaria beter
dan de $AA$ en geven hun kopieën van $a$ door. Het verlies bij de $aa$
wordt door het voordeel van de $Aa$ gecompenseerd, en $a$ blijft op 15\%.
\end{solution}

\begin{solution}{exo:g12:selection-drift-speciation:13}
De partnerkeuze op kleur: een barrière van gedrag. In troebel water faalt
die barrière en kruisen de \omterm{def:g10:biodiversity-scales:species}{soorten}, dus is er nog geen genetische
onverenigbaarheid opgebouwd: de scheiding is recent en nog omkeerbaar.
\end{solution}

\begin{solution}{exo:g12:selection-drift-speciation:14}
De hele ring rond plant elke \omterm{def:g12:selection-drift-speciation:population}{populatie} zich met haar buren voort, met
slechts kleine verschillen tussen aangrenzende; die verschillen tellen
rond de cirkel op tot de uiteinden, die elkaar ontmoeten, zich niet meer
onderling voortplanten. Het hele verloop is in één keer zichtbaar:
\omterm{prop:g12:selection-drift-speciation:speciation}{soortvorming} is een kwestie van opgetelde kleine stappen. De twee
uiteinden heten \omterm{def:g10:biodiversity-scales:species}{soorten} omdat zij zich daar waar zij elkaar ontmoeten als
\omterm{def:g10:biodiversity-scales:species}{soorten} gedragen: geen genenuitwisseling.
\end{solution}

\begin{solution}{exo:g12:selection-drift-speciation:15}
Selectie laat de dragers van sommige \omterm{def:g10:universal-dna:gene}{allelen} zich \emph{hier en nu} meer
voortplanten: beter in het nalaten van nakomelingen in déze omgeving, en
vergeleken met de andere leden van de \omterm{def:g12:selection-drift-speciation:population}{populatie} --- niet beter in enige
absolute zin. Verandert de omgeving, dan keert de richting om (de
vlinders); de drift laat \omterm{def:g12:selection-drift-speciation:population}{populaties} zonder enige reden van aanpassing
verschillen; de seksuele selectie bevoordeelt versierselen die de
overleving hinderen. De evolutie kent geen richting van verbetering,
alleen plaatselijk sorteren.
\end{solution}

\begin{solution}{pb:g12:selection-drift-speciation:1}
\textbf{1.} Kopieën van $D$: $2 \times 720 + 960 = 2400$; van $d$: $960 +
2 \times 320 = 1600$; op 4000: $p = 0.6$, $q = 0.4$.

\textbf{2.} Verwacht $0.36 \times 2000 = 720$, $0.48 \times 2000 =
960$, $0.16 \times 2000 = 320$: precies de telling. Ja.

\textbf{3.} $960/1680 \approx 57\%$.

\textbf{4.} De meeste kopieën van $d$ zitten in de donkere dragers. Na de
verwijdering: 1680 muizen, 2400 kopieën $D$, 960 kopieën $d$: $q = 960/3360 \approx
0.29$ --- een daling van 0.40, maar ver van uitroeiing.

\textbf{5.} De grote omvang: een kleine eilandpopulatie drift, dus dwalen
de \omterm{def:g12:selection-drift-speciation:population}{frequenties} door toeval (en migratie van het vasteland ontbreekt,
\omterm{def:g11:mutations:mutation}{mutatie} is verwaarloosbaar).

\textbf{6.} $DD$ 360, $Dd$ 480, $dd$ 320: 1160 dieren die zich voortplanten.

\textbf{7.} $D$: $720 + 480 = 1200$; $d$: $480 + 640 = 1120$; op 2320:
$p \approx 0.52$, $q \approx 0.48$.

\textbf{8.} Volgende generatie (per 2000): $DD$ $0.52^2 \approx 27\%$
(535), $Dd$ $\approx 50\%$ (999), $dd$ $\approx 23\%$ (466). Na de uilen:
267, 500, 466. $D$: $534 + 500 = 1034$; $d$: $500 + 932 =
1432$; $q \approx 0.58$.

\textbf{9.} Met $q = 0.58$: $DD$ 18\% (353), $Dd$ 49\% (974), $dd$ 34\%
(673); overlevenden 176, 487, 673; $d$: $487 + 1346 = 1833$ op 2672:
$q \approx 0.69$. Verloop: 0.40, 0.48, 0.58, 0.69 --- een stijging van
ongeveer 0.1 per generatie.

\textbf{10.} Elke drager van $D$ is donker en dus aan de uilen
blootgesteld, zodat de selectie bij elke generatie op elke kopie van $D$
aangrijpt; $q$ stijgt gestaag. Een \omterm{def:g11:genes-and-disease:genetic}{recessief} \omterm{def:g10:universal-dna:gene}{allel} zou zich, eenmaal
zeldzaam, in heterozygoten verschuilen, maar een \omterm{def:g11:genes-and-disease:genetic}{dominant} \omterm{def:g10:universal-dna:gene}{allel}
verschuilt zich nooit: $D$ kan tot nul worden teruggebracht.

\textbf{11.} $D$: $4 + 3 = 7$; $d$: $3 + 2 = 5$; $q = 5/12 \approx
0.42$, door het toeval van de trekking dicht bij de 0.40 van het eiland.

\textbf{12.} Onder de ouders heeft $D$ 5 kopieën en $d$ 1: $q$ in de
volgende generatie is gemiddeld $1/6$, maar de ene kopie van $d$ kan aan
geen van de jongen worden doorgegeven ($q = 0$) of aan verscheidene; het
\omterm{def:g11:enzymes-and-phenotype:phenotype}{genotype} $dd$ is al verdwenen, en met hem een groot deel van de variatie.

\textbf{13.} Met zo weinig dieren die zich voortplanten, is het bij elke
generatie een kwestie van toeval welke \omterm{def:g10:universal-dna:gene}{allelen} worden doorgegeven; een
zeldzaam \omterm{def:g10:universal-dna:gene}{allel} kan in één ongelukkig nest verloren gaan.

\textbf{14.} $d$ zal driften en binnen enkele generaties worden
vastgelegd of verdwijnen --- beide uitkomsten door toeval. Een naburig
eilandje kan vanuit hetzelfde begin met het andere \omterm{def:g10:universal-dna:gene}{allel} eindigen: de twee
eilanden verschillen in vachtkleur zonder enige reden van aanpassing.

\textbf{15.} \omterm{prop:g12:selection-drift-speciation:drift}{Genetische drift}; de piepkleine omvang van de kolonie.

\textbf{16.} Ja: zij planten zich zelden met eilandmuizen voort en de
kruisingen zijn onvruchtbaar --- voortplantingsafzondering.

\textbf{17.} De zee, een aardrijkskundige barrière; selectie in een andere
omgeving (geen uilen, andere omstandigheden) en \omterm{prop:g12:selection-drift-speciation:drift}{genetische drift} in een
kleine \omterm{def:g12:selection-drift-speciation:population}{populatie}.

\textbf{18.} Muizen die zich in verschillende seizoenen voortplanten,
ontmoeten elkaar nooit als partner: ook zonder enige aardrijkskundige
barrière stromen er geen \omterm{def:g10:universal-dna:gene}{genen}, en blijven de \omterm{def:g12:selection-drift-speciation:population}{populaties} uiteengroeien.

\textbf{19.} Nu niet: zij zijn door gedrag en onvruchtbaarheid afgezonderd
en zouden als twee \omterm{def:g10:biodiversity-scales:species}{soorten} naast elkaar blijven bestaan. Na slechts 100
jaar zouden zij zich nog vrijelijk hebben voortgeplant en weer tot één
\omterm{def:g12:selection-drift-speciation:population}{populatie} zijn samengesmolten.

\textbf{20.} $q$ van 0.40 naar ongeveer 0.69 na drie generaties uilen; op
het eilandje werd $d$ door \omterm{prop:g12:selection-drift-speciation:drift}{genetische drift} vastgelegd of verloren; twee
\omterm{def:g10:biodiversity-scales:species}{soorten} aan het eind.
\end{solution}
